\section{Topology: exact reduction and geometric obligations}

\subsection{What the exact arithmetic already proves}

Let
\[
p=12\ell_0-4\ell_1-3\ell_2.
\]
The retained exact harness verifies the determinant and Smith-normal-form calculations for the printed relation matrix. In particular, the manuscript's chosen parameters $(\ell_0,\ell_1,\ell_2)=(0,1,-1)$ give $p=-1$.

\begin{proposition}[terminal group calculation]\label{prop:terminal-group}
\Verified\ For the printed based presentation and its stated meridian relations, the resulting group is cyclic of order $|p|$. Hence the printed presentation is trivial for $(0,1,-1)$.
\end{proposition}

\begin{proof}
The exact row and column operations recorded in the audit reduce the relation matrix to Smith normal form with three unit invariant factors and final invariant factor $|p|$. A group presented by that matrix is therefore $\mathbb Z/|p|\mathbb Z$, which is trivial when $p=-1$.
\end{proof}

\begin{remark}
\Cref{prop:terminal-group} is a theorem about the printed presentation. It becomes a theorem about $\pi_1(X)$ only after the based inclusions induced by all three analytic attaching maps have been derived in the same generators and orientations.
\end{remark}

\subsection{The map that must be identified}

Write the relevant decomposition schematically as $X=U\cup N$, where $N$ is the union of the filling pieces and $U\cap N$ is the disjoint union of their collars. The homological boundary map is
\[
\delta_q=(i_{U*},-i_{N*})\colon
H_q(U\cap N;\mathbb Z)\longrightarrow
H_q(U;\mathbb Z)\oplus H_q(N;\mathbb Z).
\]
The printed matrices denoted $\alpha_q$ have verified ranks, Smith forms, and saturation properties. The central question is not another Smith-form computation; it is the equality
\[
\boxed{\alpha_q=[\delta_q]_{\mathcal B_q}}
\]
for explicitly constructed primitive oriented integral bases $\mathcal B_q$.

\begin{obligation}[finite integral model]\label{obl:finite-model}
\Open\ Construct a finite CW, simplicial, or functorial nearby-cycle model of the cusp filling and all three collars. Print the cellular maps and prove that comparison with singular homology sends them to the manuscript's bases without hidden torsion or index.
\end{obligation}

\begin{obligation}[equivariant cusp collapse]\label{obl:cusp-collapse}
\Open\ On every toric stratum, prove continuity of the collapse map, full $\bar\Lambda$-equivariance, descent to $N_0=Y_\varepsilon/\bar\Lambda$, and the asserted integral induced map. Proper discontinuity and local finiteness must be proved uniformly near the central fibre.
\end{obligation}

\begin{obligation}[based van Kampen maps]\label{obl:vankampen}
\Open\ Choose a common basepoint and paths through the collars, derive each inclusion
\[
\pi_1(U\cap N_j)\longrightarrow\pi_1(U),\qquad
\pi_1(U\cap N_j)\longrightarrow\pi_1(N_j),
\]
from the coordinate gluing, and record the translation drift, meridian orientation, primitivity, and conjugation convention. Only then may the relation matrix in \Cref{prop:terminal-group} be identified with $\pi_1(X)$.
\end{obligation}

\begin{obligation}[independent Leray route]\label{obl:leray}
\Open\ Derive, rather than normalize by convention, the global transgression
\[
d_2(12\gamma)=\pm p\,\omega
\]
in integral bases and fix the sign by comparison with an oriented local model. This route must state exactly which nearby-cycle and meridian identifications it shares with Mayer--Vietoris.
\end{obligation}

\begin{test}[fast falsification targets]
Any of the following disproves the claimed topology calculation: a nonprimitive geometric image where $\alpha_q$ assumes a primitive one; an index or torsion class omitted by specialization; a meridian sign changing $p=-1$ to a nonunit; or a Leray transgression differing from $\pm p\omega$ by a nonunit factor.
\end{test}

\subsection{Conditional recognition}

\begin{theorem}[precise conditional endpoint]
\Conditional\ Assume \Cref{obl:finite-model,obl:cusp-collapse,obl:vankampen,obl:leray} are discharged and give exactly the printed integral maps. Then the verified terminal calculations yield
\[
\pi_1(X)=0,\qquad H_k(X;\mathbb Z)=
\begin{cases}
\mathbb Z,&k=0,6,\\
0,&1\le k\le5.
\end{cases}
\]
Consequently $X$ is a smooth homotopy $6$-sphere; the standard fact $\Theta_6=0$ identifies it diffeomorphically with $S^6$.
\end{theorem}

\begin{proof}
The hypotheses identify the actual van Kampen and Mayer--Vietoris/Leray maps with the already verified presentations and matrices. Their Smith forms give the displayed fundamental group and homology. A simply connected integral homology sphere is a homotopy sphere by Hurewicz and Whitehead; the final smooth recognition uses $\Theta_6=0$.
\end{proof}
